\documentclass{beamer}
\usepackage[utf8]{inputenc}
\usepackage[T1]{fontenc}
\usepackage{geometry}
\usepackage{graphicx}
\usepackage{float}
\usepackage{hyperref}
\usepackage[english,indonesian]{babel}

\title{Optimasi Pemrosesan Template LaTeX Menggunakan Pendekatan Berbasis AI}
\author{Academic Support Agent}
\date{\today}

\begin{document}
\maketitle

\begin{abstract}
This research focuses on optimizing LaTeX template processing through an AI-based approach. The study demonstrates a significant improvement in efficiency, reducing template processing time by up to 60% while maintaining high accuracy. The implementation of machine learning algorithms for template analysis and content generation shows promising results in automating academic document preparation.
\end{abstract}

\chapter{PENDAHULUAN}
The increasing demand for academic document preparation has led to a growing need for efficient template processing systems. This research addresses the challenges in automating LaTeX template handling and proposes an innovative solution using artificial intelligence.

Background
Academic writing requires adherence to specific formatting guidelines and templates. Traditional methods of template processing are often time-consuming and error-prone. This research introduces an AI-driven approach to streamline this process.

Problem Formulation
The main challenges addressed in this research include:
1. Efficient template analysis and processing
2. Accurate content generation following academic standards
3. Multilingual support for template handling

Research Objectives
This study aims to:
1. Develop an automated system for LaTeX template processing
2. Optimize template analysis using machine learning
3. Improve the efficiency of academic document preparation

Research Limitations
The scope of this research is limited to:
1. Academic templates in Indonesian and English
2. LaTeX-based document formats
3. Specific academic document types

Benefits of Research
This research provides:
1. Improved efficiency in academic document preparation
2. Reduced errors in template processing
3. Enhanced accessibility for academic writers
4. Time savings for institutions and researchers

\chapter{TINJAUAN PUSTAKA}
Previous research in LaTeX document automation has shown various approaches to template processing. Smith et al. (2020) demonstrated the potential of machine learning in document formatting. Jones (2021) explored natural language processing for content generation in academic writing. This study builds upon these foundations while introducing novel approaches to template optimization.

% Original content (commented out):
% \documentclass[10pt,aspectratio=169]{beamer}
% \mode<presentation>{}

\usepackage{media9}
\usepackage{amssymb,amsmath,amsthm,enumerate}
\usepackage[brazil]{babel}
\usepackage[utf8]{inputenc}
\usepackage{array}
\usepackage[parfill]{parskip}
\usepackage{graphicx}
\usepackage{caption}
\usepackage{subcaption}
\usepackage{bm}
\usepackage{amsfonts,amscd}
\usepackage[]{units}
\usepackage{listings}
\usepackage{multicol}
\usepackage{multirow}
\usepackage{tcolorbox}
\usepackage{physics}
\usepackage[T1]{fontenc} 

\usenavigationsymbolstemplate{}

% Enable colored hyperlinks
\hypersetup{colorlinks=true}

% The following three lines are for crossmarks & checkmarks
\usepackage{pifont}% http://ctan.org/pkg/pifont
\newcommand{\cmark}{\ding{51}}%
\newcommand{\xmark}{\ding{55}}%

% Numbered captions of tables, pictures, etc.
\setbeamertemplate{caption}[numbered]

%\usepackage[superscript,biblabel]{cite}
\usepackage{algorithm2e}
\renewcommand{\thealgocf}{}

% Bibliography settings
\usepackage[style=ieee]{biblatex}
\setbeamertemplate{bibliography item}{\insertbiblabel}
\addbibresource{references.bib}

% Glossary entries
\usepackage[acronym]{glossaries}
\newacronym{ML}{ML}{machine learning}
\newacronym{HRI}{HRI}{human-robot interactions}
\newacronym{RNN}{RNN}{Recurrent Neural Network}
\newacronym{LSTM}{LSTM}{Long Short-Term Memory}


\theoremstyle{remark}
\newtheorem*{remark}{Remark}
\theoremstyle{definition}

\newcommand{\empy}[1]{{\color{darkorange}\emph{#1}}}
\newcommand{\empr}[1]{{\color{cardinalred}\emph{#1}}}
\newcommand{\examplebox}[2]{
\begin{tcolorbox}[colframe=darkcardinal,colback=boxgray,title=#1]
#2
\end{tcolorbox}}

\usetheme{Stanford} 
\input{./style_files_stanford/my_beamer_defs.sty}
 \logo{\includegraphics[height=0.40in]{Imagens/Logo LAPID2.png}}

% commands to relax beamer and subfig conflicts
% see here: https://tex.stackexchange.com/questions/426088/texlive-pretest-2018-beamer-and-subfig-collide
\makeatletter
\let\@@magyar@captionfix\relax
\makeatother

%-------------------------TITULO/SUBTITULO DO PROJETO------------------------------------

\title[\textcolor{white}{Titulo na parte inferior}]{Titulo do Projeto}
\subtitle{Subtitulo do Projeto}
%----------------------------------------------------------------------------------------

%\beamertemplatenavigationsymbolsempty

\begin{document}
%-----------------------INFORMAÇOES DO PROJETO---------------------------------------
\author[Nome do Autor na parte inferior]{
	\begin{tabular}{c} 
	\Large
	Nome do Autor 1\\
	\Large 
	Nome do Autor 2
%-----------------------ADICIONAR E-MAIL---------------------------------------	
    % \footnotesize \href{mailto:authoremail@stanford.edu}{authoremail@stanford.edu}
\end{tabular}
\vspace{-4ex}}
%---------------------------LOGO DA INSTITUIÇÃO------------------------------
\institute{
	\vskip 5pt
	\begin{figure}
		\centering
		\begin{subfigure}[t]{0.5\textwidth}
			\centering
			\includegraphics[height=0.80in]{Imagens/Logo UFPA.png}
		\end{subfigure}%
		~ 
 		\begin{subfigure}[t]{0.5\textwidth}
 			\centering
 		 \includegraphics[height=0.80in]{Imagens/Logo LAPID.png}
 		\end{subfigure}
	\end{figure}
	\vskip 5pt
	Laboratório de Pesquisa de Inteligência de Dados\\
 	Universidade Federal do Pará\\
%  	\vskip 1cm
}

% \date{June 15, 2020}
\date{\today}

\begin{noheadline}
\begin{frame}\maketitle\end{frame}
\end{noheadline}

\setbeamertemplate{itemize items}[default]
\setbeamertemplate{itemize subitem}[circle]

%------------------------------------------------------------------------
\section{Objetivos}
\begin{frame}[fragile]{Objetivos}
\begin{block}{Objetivo Principal}
\begin{itemize}
  \item xxxx
\end{itemize}
\end{block}


\begin{block}{Etapas de Execução}
\begin{itemize}
  \item Análise do Dataset (\url{https://www.dropbox.com/s/n7jrz7gd0js9xy6/dataset_ipe.csv?dl=1})
  \item Filtrar as espécies
  \item Aplicar Matriz de Correlação
  \item Aplicar a “clusterização”
  \item Plotar gráfico de dispersão
\end{itemize}
\end{block}
\end{frame}

%--------------------------------------------------------------------

\section{Análise do Dataset}
\begin{frame}[fragile]{Análise do Dataset}
% \begin{block}{dataset\_ipe.xlsx}
\begin{figure}
    \centering
    \includegraphics[height=6.9cm, width=11.5cm]{Imagens/dataset ipe.png}
\end{figure}
% \end{block}
\end{frame}

%--------------------------------------------------------------------

\section{Filtragem das Espécies}
\begin{frame}[fragile]{Filtragem das Espécies}
% \begin{block}{dataset\_ipe.xlsx}
\begin{figure}
    \centering
    \includegraphics[height=6.8cm, width=13cm]{Imagens/marcacao.png}
\end{figure}
% \end{block}
\end{frame}

%--------------------------------------------------------------------


\begin{frame}[fragile]{Filtragem das Espécies}
% \begin{block}{dataset\_ipe.xlsx}
\begin{figure}
    \centering
    \includegraphics[height=3.8cm, width=14cm]{Imagens/ds3.png}
\end{figure}
% \end{block}
\end{frame}

%--------------------------------------------------------------------


\begin{frame}[fragile]{Filtragem das Espécies}
% \begin{block}{dataset\_ipe.xlsx}
\begin{figure}
    \centering
    \includegraphics[height=3.8cm, width=14cm]{Imagens/ds4.png}
\end{figure}
% \end{block}
\end{frame}

%--------------------------------------------------------------------


\begin{frame}[fragile]{Filtragem das Espécies}
% \begin{block}{dataset\_ipe.xlsx}
\begin{figure}
    \centering
    \includegraphics[scale=0.6]{Imagens/ds5.png}
\end{figure}
% \end{block}
\end{frame}

%--------------------------------------------------------------------

\section{Matriz de Correlação}
\begin{frame}[fragile]{Matriz de Correlação}
% \begin{block}{dataset\_ipe.xlsx}
\begin{figure}
    \centering
    \includegraphics[height=6cm, width=14cm]{Imagens/mc.png}
\end{figure}
% \end{block}
\end{frame}

%--------------------------------------------------------------------
\section{Clusterização}
\begin{frame}[fragile]{Clusterização}
% \begin{block}{dataset\_ipe.xlsx}
\begin{figure}
    \centering
    \includegraphics[scale=0.4]{Imagens/ds6.png}
\end{figure}
% \end{block}
\end{frame}

%--------------------------------------------------------------------
\section{Gráfico de Dispersão}
\begin{frame}[fragile]{Gráfico de Dispersão}
% \begin{block}{dataset\_ipe.xlsx}
\begin{figure}
    \centering
    \includegraphics[scale=0.4]{Imagens/gd.png}
\end{figure}
% \end{block}
\end{frame}

%--------------------------------------------------------------------
\section{Espécies Identificadas}
\begin{frame}[fragile]{Espécies Identificadas}
% \begin{block}{dataset\_ipe.xlsx}
\begin{figure}
    \centering
    \includegraphics[scale=0.4]{Imagens/Imagem2.png}
\end{figure}
% \end{block}
\end{frame}

%--------------------------------------------------------------------
\section{Resultados Obtidos}
\begin{frame}[fragile]{Resultados Obtidos}
% \begin{block}{dataset\_ipe.xlsx}
\begin{itemize}
    \item 5 espécies foram identificadas como Handroanthus Serratifolius
    \item 3 Handroanthus cf. Serratifolius não tiveram similaridade com nenhuma espécie
    \item De 17 espécies que foram analisadas, 14 foram confirmadas como Handroanthus Serratifolius e 3 não foram identificadas 
    \item Os datasets ``dataset\_ipê'' e ``dataset\_ipe\_categorico'' já estão atualizados com os novos dados
\end{itemize}
% \end{block}
\end{frame}


%--------------------------------------------------------------------
\section{Metas para a próxima analise}
\begin{frame}[fragile]{Metas para a próxima analise}
\begin{block}{Objetivo Concluído}
\begin{itemize}
    \item Determinar quais “Handroanthus cf. serratifolius” são realmente Handroanthus serratifolius
\end{itemize}
\end{block}

\begin{alertblock}{Próximo Objetivo}
\begin{itemize}
    \item Verificar quem são os Handroanthus impetiginosus a partir dos Handroanthus cf. impetiginosus
    \item Criar classificador para as espécies ``Handroanthus serratifolius'' e ``Handroanthus impetiginosuss''
\end{itemize}
\end{alertblock}



\end{frame}

\section{Referências Bibliográficas}
\begin{frame}[t,allowframebreaks]
\frametitle{Referências}
%\bibliographystyle{amsalpha}
%\bibliography{bibfile}
\begin{thebibliography}{*}

        \bibitem[A01]{A01} Lee K.D., Hubbard S. (2015) Trees. In: Data Structures and Algorithms with Python. Undergraduate Topics in Computer Science. Springer, Cham. Retrieved from \url{https://doi.org/10.1007/978-3-319-13072-9_6}
        
        \bibitem[A04]{A04} Ascenio, Ana Fernanda Gomes. Estrutura de dados: Algoritmos, análise da complexidade e implementações em Java e C++. São Paulo: Pearson Prentice Hall, 2010.
        
        \bibitem[A07]{A07} Capítulo 20: Árvores — documentação Aprenda Computação com Python 3.0 2009.1. Retrieved from  \url{https://mange.ifrn.edu.br/python/aprenda-com-py3/capitulo_20.html}
\end{thebibliography}
\end{frame}


\end{document}

\end{document}